\section{Discussion}

The results of the synthesis and verification phases reveal several key architectural insights regarding the acceleration of the 2D DCT on an FPGA fabric.

\subsection{Architectural Tradeoffs: Pipelining vs. Parallelism}
The progression from D1 to D4 empirically demonstrates the highly nuanced impacts of pipelining and parallelism. Pipelining (D2) proved to be the ``sweet spot'' for strict area constraints. Because ALMs naturally contain unused registers, D2 nearly doubled the $F_{max}$ and throughput over the D1 baseline while incurring a negligible 130-ALM penalty, making it a highly efficient upgrade. Conversely, spatial parallelism without pipelining (D3) emerged as a potential design trap. While it successfully reduced clock cycle latency by a factor of 8, the massive combinatorial adder trees significantly increased the critical path, dropping $F_{max}$ to 85 MHz. As a result, D3's overall throughput-per-ALM efficiency decreased below even the D1 baseline. This indicates that if a platform can afford the area footprint of D3, the designer should absolutely step up to D4. The combined D4 design proves that these techniques are orthogonal and must be layered: by pipelining the parallel arrays, D4 achieved both the highest absolute throughput (3.04M blocks/sec) and the highest overall throughput-per-ALM efficiency, making it the undisputed optimal choice for high-performance applications.

\subsection{Memory Bottlenecks}
A significant observation from the hardware breakdown is the outsized resource consumption of the transpose buffer memory. Because the row-column decomposition inherently requires the full $8 \times 8$ intermediate matrix to be stored before column processing can begin, the transpose buffer acts as an unavoidable physical bottleneck. In the D1 and D2 designs, the memory bits drastically overshadow the core computational logic. While unavoidable in a 1D-decoupled architecture, future iterations could explore custom dual-port block RAM instantiations rather than inferred distributed RAM to minimize the routing overhead associated with this memory structure.

\subsection{Scalability for HD and 4K/8K Video}
The modular, parameterized nature of this SystemVerilog implementation ensures that the core engine is highly scalable. Processing video streams in real-time requires immense throughput; a single $1920 \times 1080$ HD frame contains 32,400 $8 \times 8$ pixel blocks. At standard 60 frames per second (FPS), this demands a throughput of $1.94$ million blocks/sec. The D4 variant, achieving 3.04 million blocks/sec, easily handles a 60 FPS HD stream single-handedly. However, scaling to 4K ($3840 \times 2160$) at 60 FPS requires $7.78$ million blocks/sec, and 8K ($7680 \times 4320$) at 60 FPS demands a staggering $31.1$ million blocks/sec. To meet these ultra-high-definition targets, the top-level parameterization allows multiple D4 engines to be instantiated in parallel. Specifically, a 4K 60 FPS stream would require instantiating 3 parallel D4 cores, while an 8K 60 FPS stream would require 11 cores. This multi-core distribution fabric scales performance linearly without requiring any fundamental redesign of the core MAC array.

\subsection{Fixed-Point Limitations}
While the 16-bit Q2.14 fixed-point formatting achieved a highly respectable 51.05 dB PSNR, there remains an inherent hard ceiling on the reconstruction fidelity. The strict $\pm 2$ LSB quantization error is mathematically unavoidable without widening the datapath. If a specific medical or scientific imaging application requires mathematically lossless reconstruction, the hardware would need to be re-parameterized to support 32-bit arithmetic. However, this would heavily penalize the $F_{max}$ and effectively double the size of the required transpose buffer.
